\documentclass[11pt]{article}

% --- arXiv-friendly preamble (standard packages only) ---
\usepackage[utf8]{inputenc}
\usepackage[T1]{fontenc}
\usepackage{amsmath,amssymb,amsthm}
\usepackage{booktabs}
\usepackage{geometry}
\usepackage{hyperref}
\usepackage{xcolor}
\geometry{margin=1in}
\hypersetup{colorlinks=true,linkcolor=blue,citecolor=blue,urlcolor=blue}

\newtheorem{definition}{Definition}
\newtheorem{proposition}{Proposition}
\newtheorem{theorem}{Theorem}

\newcommand{\dcausal}{d_{\mathrm{c}}}
\newcommand{\Nk}{\mathcal{N}_k}
\newcommand{\code}[1]{\texttt{#1}}

\title{\textbf{The Cognitive MMU:\\ Deterministic, Auditable Working Memory for LLM Coding Agents\\ \large (and an honest account of why causal retrieval does not beat RAG)}}
\author{The CCOS Project\thanks{Causal Context Operating System (CCOS), an open
research prototype. Source, reproduction scripts and this paper:
\url{https://github.com/CHECKUPAUTO/CCOS}.}}
\date{\today}

\begin{document}
\maketitle

\begin{abstract}
An LLM coding agent's working memory is managed opaquely: a retrieval stack (RAG,
GraphRAG, agent frameworks) pages a repository into a bounded window and mutates
that window probabilistically as a session unfolds, so when an agent drifts after
dozens of steps it is impossible to say \emph{why} or \emph{when} its
representation of the project corrupted. CCOS treats agent memory the way an
operating system treats RAM: source is parsed into an \emph{event-sourced} causal
graph, nodes are scored and paged under a token budget, and every state transition
is recorded in a hash-chained log that replays bit-for-bit. We give a precise,
deterministic definition of a \emph{causal region} (a connected component of the
cross-file causal-link graph) and of causal distance (a weighted shortest path),
and prove that regions and the paged window are a pure function of the graph,
reconstructing exactly on replay --- the property that makes an agent's memory
\emph{auditable} and \emph{time-travel debuggable}. We then ask whether this causal structure makes CCOS a \emph{better retriever}
than a baseline, and report --- honestly --- that it does not. Measured on $70$
real bug-fix commits across three mature crates, causal selection \textbf{ties a
plain lexical TF-IDF retriever} and loses to it at a tight budget; on real bugs a
fix's files share vocabulary, so lexical similarity finds them too, and a
proof-of-concept that seeds CCOS from a crash trace is likewise beaten by
RAG-over-the-error-message (Rust error messages name the cause). An LLM-free
simulation \emph{does} separate the two when a cause is constructed to be
lexically dissimilar from its symptom, and a first on-device run with a $7$B model
reproduces it --- but that regime does not occur on real code. The contribution is
therefore \emph{not} better retrieval. It is the infrastructure a retrieval stack
structurally lacks: a \emph{deterministic, replayable, auditable} agent memory in
which one can rewind the exact context state at any step and replay it under
different parameters (\emph{time-travel debugging}). To our knowledge CCOS is the
first system to treat an agent's \emph{working memory itself} as a transactional
subsystem --- deterministic, hash-chained, replayable bit-for-bit, and
\emph{post-mortem debuggable}: a flight recorder for attention, with a breakpoint
on the exact step at which the right information was evicted from the window. We release the engine, the
honest validation harness, and this paper, labelling throughout what is
\emph{proven} (formal structure, determinism, replay), \emph{measured} (retrieval
coverage --- where CCOS does \emph{not} beat lexical RAG on real bugs), and the one
axis where CCOS \emph{does} win: \textbf{efficiency}. Run end-to-end on real
single-file fixes with a 30B model and a compiler-in-the-loop retry, CCOS and RAG
resolve \emph{equally} ($2/10$), but CCOS does so on $776$ context tokens against
the lexical retriever's $5366$ --- a $6.9\times$ reduction (and $4.1$--$9.1\times$
across $51$ fixes from three crates, model-free), because the causal region stops
at the working set rather than padding a top-$k$ to budget. CCOS is thus not a
better retriever or solver, but a more \emph{frugal}, \emph{auditable} one --- it
self-calibrates, with no $k$ to tune. (On its own source tree, region selection covers a task's $k$-hop
neighbourhood with recall $0.97$ vs $0.35$ flat at $\approx\!48\%$ fewer tokens;
regions are $95.5\%$ internally connected.)
\end{abstract}

\section{Introduction}

An LLM coding agent operating over a repository must repeatedly decide \emph{what
to put in the context window}. The dominant pattern is retrieval: embed code
chunks, rank them against the task, and concatenate the top-$k$ until a token
budget is exhausted~\cite{lewis2020rag}. This treats context as a one-dimensional
ranked list. Two failure modes follow. First, \emph{dilution}: as tasks lengthen,
the window accumulates globally salient but task-irrelevant code, and models
attend poorly to the middle of long contexts~\cite{liu2023lost}. Second,
\emph{incoherence}: the top-$k$ chunks need not form a connected unit of
reasoning --- a function may be paged in without its callers, its error sites or
the data it depends on.

Operating systems faced the analogous problem decades ago and answered with the
\emph{working set}: the set of recently referenced pages that must stay
resident~\cite{denning1968working}. CCOS adopts the analogy for LLM
context~\cite{packer2023memgpt}: source is parsed into a causal
graph, nodes are scored, and a bounded ``context window'' is paged in and out
like RAM~$\leftrightarrow$~VRAM. Unlike a retrieval stack, every transition is
recorded in a hash-chained event log, so the memory is not a black box but an
\emph{auditable, replayable} artefact. Our contributions are:

\begin{enumerate}
  \item A \textbf{formal, deterministic model of context regions}
  (\S\ref{sec:regions}, \S\ref{sec:determinism}): causal distance as a weighted
  shortest path, region membership as a connected component of the cross-file
  causal-link graph, and a determinism theorem --- regions and the paged window
  are a pure function of the graph, so a session reconstructs bit-for-bit from its
  hash-chained event log (tamper-evident replay).
  \item \textbf{Event-sourced agent sessions with time-travel debugging}
  (\S\ref{sec:timetravel}): every cognitive operation (ingest, failure signal,
  recall) is logged; the exact context state at any step is reconstructible; and a
  recall can be \emph{replayed under different parameters} to ask whether the agent
  would have decided better --- the capability a probabilistic retrieval stack
  structurally lacks. To our knowledge this is the first treatment of context
  assembly \emph{itself} as a replayable, post-mortem-debuggable subsystem --- a
  \emph{flight recorder for an agent's attention} --- with an \emph{eviction
  watchpoint} that names the exact step and operation at which the true cause was
  squeezed out of the budgeted window.
  \item An \textbf{honest validation harness and a negative result}
  (\S\ref{sec:protocol}): on $70$ real bug-fix commits, causal selection
  \emph{does not} beat a lexical TF-IDF retriever at placing a fix's files in the
  window (it ties, and loses at a tight budget), and a crash-trace pivot is beaten
  by RAG-over-the-error-message. We report this plainly --- it relocates CCOS's
  value from \emph{retrieval} to \emph{auditability}.
  \item An \textbf{LLM-free locality measurement} (\S\ref{sec:eval}) with real,
  reproducible numbers, and a falsifiable \emph{sufficient}-condition protocol
  (Phase~4) we specify but leave open.
\end{enumerate}

\section{Related Work}\label{sec:related}

\paragraph{Retrieval-augmented generation.} RAG augments a parametric model with
a non-parametric store and retrieves passages per
query~\cite{lewis2020rag}. Self-RAG adds reflection tokens that gate
retrieval and critique generations~\cite{asai2023selfrag}. These operate on
\emph{independent} chunks; coherence between retrieved items is not modelled.

\paragraph{Graph- and structure-aware retrieval.} GraphRAG builds an entity
knowledge graph and answers global queries by summarising community
structure~\cite{edge2024graphrag}. For code, property graphs unify AST, control
flow and data flow~\cite{yamaguchi2014cpg}. CCOS regions are in this lineage but
target \emph{paging}: which connected sub-structure to make resident for a task,
under a token budget, with deterministic eviction.

\paragraph{Agent memory.} MemGPT casts the LLM as an OS that pages between an
in-context window and external storage~\cite{packer2023memgpt}; Generative Agents
retrieve memories scored by recency, importance and
relevance~\cite{park2023generative} --- the same factors CCOS aggregates into a
region temperature. LangGraph~\cite{langgraph} structures agents as stateful
graphs of steps; it orchestrates \emph{control} flow, whereas CCOS structures
\emph{memory}. The two are complementary.

\paragraph{Context-window management.} KV-cache eviction keeps ``heavy hitter''
or sink tokens resident~\cite{liu2023lost}, and PagedAttention applies OS paging
to the KV cache~\cite{kwon2023paged}. These act at the token level inside one
forward pass; CCOS acts at the \emph{semantic} level across an agent session.

\section{The Causal Context Substrate}\label{sec:substrate}

CCOS parses source into a directed \emph{causal memory graph} $G=(V,E)$. A node
$v\in V$ is a file, module, symbol or external dependency, with scalar fields
$\mathrm{imp}(v)$ (base importance), $\mathrm{fail}(v)\in[0,1]$ (failure
relevance) and $\mathrm{rec}(v)\in[0,1]$ (recency). An edge $e=(u\!\to\!w)\in E$
carries a weight $w(e)\in(0,1]$ and a type (containment, dependency, reference,
causation). The kernel assigns each node a causal score
\begin{equation}
\mathrm{score}(v) = \mathrm{clamp}\big(0.15\,\mathrm{imp}(v) + 0.50\,\mathrm{fail}(v)
  + 0.30\,\mathrm{rec}(v) + 0.05\ln(1{+}\mathrm{acc}(v)),\,0,1\big),
\label{eq:score}
\end{equation}
where $\mathrm{acc}(v)$ is the access count. Faults propagate along edges,
$\mathrm{fail}$ decays with a logical clock, and every state transition is
appended to a hash-chained event log enabling deterministic
replay~\cite{haber1991timestamp}. Node identifiers are namespaced
(\code{file:p}, \code{mod:p:n}, \code{use:p:path}, \code{sym:p:n},
\code{dep:root}); the owning file of a node is recoverable from its identifier.
CCOS v0.2 pages in the top-$\mathrm{score}$ nodes: a flat, 1-D policy. We now
make selection spatial.

\section{Context Regions}\label{sec:regions}

\subsection{Causal distance}

\begin{definition}[Causal distance]
Let $\hat G$ be the undirected multigraph on $V$ induced by $E$, and assign each
edge the cost $c(e) = -\ln w(e) \ge 0$, so that a stronger causal link is a
shorter step. The \emph{causal distance} $\dcausal(u,v)$ is the minimum total
cost over all $u$--$v$ paths in $\hat G$, and $+\infty$ if none exists. The
unweighted hop distance $\mathrm{hops}(u,v)$ is defined analogously with unit
costs.
\end{definition}

$c(e)\ge 0$ since $w(e)\le 1$, so $\dcausal$ is a genuine shortest-path metric on
each connected component (non-negative, symmetric, triangle inequality). The
\emph{$k$-hop causal neighbourhood} of a target $t$ is
\begin{equation}
\Nk(t) = \{\, v \in V : \mathrm{hops}(t,v) \le k \,\}.
\end{equation}
$\Nk(t)$ is the ground truth ``what is causally relevant to a task at $t$'' used
in the evaluation (\S\ref{sec:eval}).

\subsection{Region membership}

External dependency hubs (e.g.\ \code{dep:std}) connect almost everything and
must not collapse the graph into one component. We therefore separate them.

\begin{definition}[Cross-file causal link]
For a non-external node $v$ let $\phi(v)$ be its owning file. Two files $f,g$ are
\emph{directly linked}, written $f \approx g$, iff there exists an edge
$(u\!\to\!w)\in E$ with $\phi(u)=f$, $\phi(w)=g$, $f\neq g$, and neither $u$ nor
$w$ is an external dependency node.
\end{definition}

\begin{definition}[Region]\label{def:region}
Let $\approx^{*}$ be the reflexive-transitive closure of $\approx$ on the set of
file keys. A \emph{region} is the set of all nodes whose files lie in one
equivalence class of $\approx^{*}$; all external dependency nodes form one
additional region. Two nodes belong to the same region iff their files are in the
same connected component of the cross-file causal-link graph.
\end{definition}

\begin{proposition}[Regions partition the nodes]
The regions of Definition~\ref{def:region} form a partition of $V$: every node
lies in exactly one region.
\end{proposition}
\begin{proof}
$\approx^{*}$ is an equivalence relation on file keys, so its classes partition
the file keys; the external bucket is a distinct class by construction. Mapping
each non-external node to its file's class and each external node to the external
class is a total function into disjoint blocks, hence a partition. \qed
\end{proof}

By default (only containment and import edges) each source file is its own region.
A genuine cross-file dependency or a propagated failure merges files into a single
multi-file region --- the ``zone of knowledge'' an agent should wake together.

\subsection{Region scalars}

For a region $R$ with member set $M$:
\begin{align}
\mathrm{heat}(v)        &= 0.5\,\mathrm{score}(v) + 0.3\,\mathrm{fail}(v) + 0.2\,\mathrm{rec}(v), \\
\mathrm{temp}(R)        &= \mathrm{clamp}\!\Big(\tfrac{1}{|M|}\textstyle\sum_{v\in M}\mathrm{heat}(v),\,0,1\Big), \label{eq:temp}\\
\mathrm{dens}(R)        &= \frac{|\{\,e\in E : \text{both endpoints} \in M\,\}|}{|M|}. \label{eq:dens}
\end{align}
Temperature is how ``awake'' a region is; density is its internal causal cohesion
(internal edges per member). Activation warms a region
($\mathrm{temp}\mathrel{+}=0.25$, capped) and records a logical tick; a cooldown
step multiplies temperatures by a decay factor and evicts regions below a floor.

\subsection{Dynamic admission policy}\label{sec:policy}

The static threshold $0.6$ becomes a function of token pressure $u\in[0,1]$ (the
used fraction of the budget) and task complexity $\kappa\in[0,1]$:
\begin{align}
\theta(u,\kappa)       &= \mathrm{clamp}(0.6 + 0.3\,u - 0.2\,\kappa,\; 0.05,\; 0.95), \\
a(R)                   &= 0.55\,\mathrm{temp}(R) + 0.30\,\frac{\mathrm{dens}(R)}{1+\mathrm{dens}(R)} + 0.15\,\kappa, \\
\mathrm{admit}(R)      &\iff a(R) \ge \theta(u,\kappa).
\end{align}
A hot, cohesive region can be admitted even when the static $0.6$ would reject it;
a nearly-full window raises $\theta$ so only the hottest regions enter.

\section{Determinism and Replay}\label{sec:determinism}

\begin{theorem}[Regional determinism]
The region partition, every region scalar in
\eqref{eq:temp}--\eqref{eq:dens}, and the activation history are pure,
order-independent functions of the graph $G$ and the sequence of (logically
timestamped) activation events. Consequently, given the event log of a session,
the engine state reconstructs identically: if $G'$ is the graph rebuilt from the
log and $L$ its region events, then $\mathrm{replay\_from}(G',L)$ equals the live
engine that produced $L$.
\end{theorem}
\begin{proof}[Proof sketch]
Clustering enumerates nodes and edges in sorted order and computes connected
components by a sorted breadth-first search, so its output is independent of hash
iteration order. Equations \eqref{eq:score}, \eqref{eq:temp} and \eqref{eq:dens}
are deterministic arithmetic over node fields and a count of qualifying edges.
Activation reads a logical clock (a counter), never wall-clock time, and the
emitted \code{RegionActivated} event records the exact tick. Reconstruction
re-clusters $G'$ (identical base state, since $G'\!=\!G$ structurally) and applies
the recorded activations in log order; each step is the same pure function of the
same inputs, so the resulting state is identical. The integration test
\code{replay\_reconstructs\_identical\_engine} checks
$\mathrm{engine}=\mathrm{replay\_from}(G',L)$ by structural equality. \qed
\end{proof}

This extends CCOS's existing guarantees: the primary event log is hash-chained and
tamper-evident, so the region history is auditable, and a 10\,000-cycle test
exhibits no drift in region count or temperatures.

\paragraph{An auditable input boundary.} The same deterministic, replayable
substrate hardens the text an agent ingests. Hidden-character injection vectors
--- bidirectional overrides (the \emph{Trojan Source} attack, CVE-2021-42574),
zero-width formatting, and the Unicode \emph{Tags} block used for invisible ASCII
smuggling --- are de-obfuscated at ingest into explicit, visible literals
(\code{[U+202E RLO]}) rather than silently stripped, and the findings are
recorded against the same hash-chained log, so a replay reproduces the cleaned
state. A downstream linear log-space classifier --- the closed form of
multinomial Naive Bayes, $\mathrm{logit}=b+W\!\cdot\!X$ over a hashing-trick
feature vector, its weights locked in a checksum-verified blob --- adds a
\emph{deterministic, forensically decomposable} injection signal (a held-out
red-team measures $F_1=0.90$; precision $0.87$, recall $0.93$). We are deliberate
that this is a \emph{signal, not a solution}: a bag-of-features model is evaded by
paraphrase, and no character-level pass addresses semantic injection --- the real
mitigation is privilege separation in the host. The contribution is not a novel
detector but that de-obfuscation and scoring inherit CCOS's defining property:
every security decision is reproducible and auditable down to the exact feature
that moved it.

\section{Event-sourced sessions and time-travel debugging}\label{sec:timetravel}

Determinism is not only a correctness property; it is the basis of the capability
that, our evaluation will argue, distinguishes CCOS from a retrieval stack. An
\code{AgentSession} records every cognitive operation an agent performs on its
memory --- \code{Ingest}, \code{SignalFailure}, \code{Recall} --- as an ordered
timeline. Because each operation is a deterministic function of the prior state,
the memory after the first $n$ operations is reconstructed exactly by replaying
the mutating ones (\code{replay\_to(n)}): one can \emph{rewind the agent's mind} to
any step.

The operation that matters for debugging is the counterfactual. \code{recall\_what\_if($n$,
$q$, $b$)} replays the memory to step $n$ and re-runs a recall with a different
query $q$ or token budget $b$, returning the window the agent \emph{would} have
seen. When an agent emits a bad patch at step 15, an operator can replay its exact
context at step 14, widen the budget or change the anchor, and observe whether the
decision improves --- a closed-loop, reproducible debugger for an agent's working
memory.

\paragraph{A breakpoint on attention.} Replay also localises \emph{drift}. The
\code{missing} watchpoint scans the timeline for the exact step at which a given
node --- typically the true cause of a failure --- was squeezed out of the budgeted
window by competing pressure, naming the triggering operation and the token gap;
a complementary \code{energy} view surfaces the causal-heat migration through the
graph that a file-level diff misses. It is, in effect, a \emph{breakpoint on an
agent's attention}: the precise moment the right information left the window. An
opaque retrieval stack can show only the final, corrupted context; CCOS can point
to the step --- and the operation --- at which it went wrong. A RAG or framework
stack mutates its store probabilistically across a
session and keeps no canonical, replayable transition log, so the same question
(\emph{why, and when, did the representation corrupt?}) has no answer there. We do
not claim this improves task success; we claim it is a property the alternatives do
not have, and that it is the honest locus of CCOS's contribution
(\S\ref{sec:protocol}).

\paragraph{The log as training data.} Because the timeline is deterministic and
replayable, it is also a \emph{counterfactual evaluation set} for retrieval
itself. We read a reward straight off it: for each recorded recall, was the node
the agent engaged \emph{next} --- a failure signal or page-fault, taken as a
post-hoc relevance signal --- present in the window that recall \emph{would} have
produced under candidate scoring weights? Maximising this hit rate by deterministic
coordinate ascent (each candidate scored by a full replay) yields scoring weights
tuned to the session's own history; adopting them is recorded as an operation, so
the learned policy is itself auditable and reproduced on replay. The reward is an
honest proxy and the optimiser is greedy, so we claim only that this is retrieval
that \emph{trains on a substrate the alternatives lack}: a probabilistic RAG store
keeps no canonical, replayable log to learn from. This closes the loop the title
implies --- the working memory is not only inspectable but \emph{self-improving}
from its own auditable record.

\section{Implementation}\label{sec:impl}

The engine is $\approx\!600$ lines of dependency-light Rust layered above the
graph, with no change to the parser, guard, incremental builder or event-sourcing
core. The public surface is: \code{ContextRegionEngine} (\code{cluster\_nodes},
\code{initialize\_regions}, \code{activate\_region}, \code{tick\_cooldown},
\code{replay\_from}), the data types \code{ContextRegion} / \code{ContextPoint} /
\code{ContextWindow}, the \code{ContextPolicy}, and five new event variants
(\code{RegionCreated}, \code{RegionActivated}, \code{RegionMerged},
\code{RegionEvicted}, \code{ContextWindowGenerated}). A \code{ccos regions} CLI
exposes clustering, activation and the locality report. The whole crate builds
warning-clean under \code{clippy -D warnings} and passes 431 tests.

\paragraph{The cognitive MMU, made real.} Paging lives in the code, not only the
prose: eviction is \emph{non-destructive}. When the resident node set exceeds its
cap, the lowest-scored nodes are \emph{demoted}---with their incident edges---into
a COLD tier rather than dropped, and any node faults back on demand
(\code{page\_in}). A failure signal, or a recall \emph{around} a demoted node,
resurrects it and its cold neighbours transparently; the page-in is a
\emph{replayed} side effect, so \mbox{\code{replay}\,$=$\,\code{live}} still holds.
The COLD tier's unbounded \emph{content} optionally spills to a content-addressed
on-disk store (SHA-256 keys, deduplicated, and \emph{hash-verified} on read---a
tampered or missing blob is a cold-miss, never a silent empty restore), bounding
the resident-plus-cold \emph{content} footprint in RAM while the backing store on
disk is unbounded. At the deepest tier, once the backing store \emph{itself} must
stay frugal, an optional budget \emph{compacts} the coldest content---code
skeletonised, prose summarised---to a causal summary, discarding the original. This
is genuine, \emph{auditable forgetting}: it is lossy, but observable
(\code{cold\_compacted}) and never a silent drop. This is the sense in which working
memory is ``unbounded'': a \emph{direction}, frugality~$\times$~available RAM, rather
than a literal claim---and at the floor frugality is the master, exactly as the
demand-paging analogy demands. Honestly scoped: spill moves only \emph{content} to
disk (per-cold-node metadata still grows in RAM---an on-disk index is future work),
blobs are stored verbatim (dedup, no compression codec yet), and compaction bounds
the cold \emph{content} footprint, not the entry \emph{count}. Both spill and
compaction are opt-in; the default path keeps them off and serialization
byte-identical, so the determinism guarantees of \S\ref{sec:determinism} are
untouched.

\paragraph{Hybrid entry fusion.} A free-text recall resolves its entry node by
\emph{reciprocal-rank fusion}~\cite{cormack2009reciprocal} of three independent
rankings---lexical token overlap, semantic INT4-TF-IDF cosine, and a causal
\emph{active-failure} focus---before the usual region expansion. RRF compares
ranks, not raw scores, so the three incomparable signals fuse without
calibration. The causal vote is deliberately \emph{sparse}: it ranks only nodes
under failure pressure, so it abstains on a quiet graph (no spurious id-ordered
bias when scores are flat) and speaks for the active problem region once a failure
is signalled---the CCOS-native attention signal rather than a generic importance
prior. Deterministic, and an honest improvement to \emph{entry selection} only:
the downstream region expansion and the evaluation of \S\ref{sec:eval} are
unchanged. The semantic signal itself is, by default, deterministic INT4 TF-IDF;
an opt-in build distils it into a learned \emph{latent-semantic} projection (LSA /
truncated SVD of the corpus's document--term matrix, via a fixed Jacobi sweep),
which captures synonymy raw TF-IDF cannot while remaining zero-dependency and
deterministic --- the ``distil, don't couple'' rule that lets a learned component
respect the replay invariant.

\section{Evaluation: what we can measure today}\label{sec:eval}

We separate \emph{measured} results (this section, fully reproducible and
LLM-free) from \emph{hypothesised} agent-level gains (\S\ref{sec:protocol}).

\paragraph{Setup.} We run on CCOS's own source tree
($|V|=705$ nodes, $|E|=822$ edges, yielding $26$ regions). For a target node $t$
we take $\Nk(t)$ with $k=2$ as ground truth and compare two selection strategies
at the region's size budget: \textbf{flat} --- the top-$|R|$ nodes by
$\mathrm{score}$ (CCOS v0.2 / classical ranked retrieval) --- and \textbf{region}
--- the members of $t$'s region. We report causal precision $|S\cap\Nk|/|S|$,
recall $|S\cap\Nk|/|\Nk|$, and the tokens each strategy needs to \emph{cover}
$\Nk$. All numbers are emitted by \code{scripts/region\_benchmark.sh} and are
deterministic.

\begin{table}[h]
\centering
\begin{tabular}{lcc}
\toprule
Strategy & causal precision & causal recall \\
\midrule
flat (v0.2 / ranked) & 0.021 & 0.347 \\
\textbf{region (v0.3)} & \textbf{0.057} & \textbf{0.972} \\
\bottomrule
\end{tabular}
\caption{Locality on \code{src/} ($k=2$, 12 targets). Region selection covers
$97\%$ of a task's causal neighbourhood versus $35\%$ for flat, at
$\approx\!48\%$ fewer tokens for equal coverage.}
\label{tab:locality}
\end{table}

\paragraph{Cohesion and cost.} Regions reach an \textbf{average causal density of
$0.955$} (Eq.~\ref{eq:dens}, normalised) --- they are almost entirely internally
connected, i.e.\ genuine causal clusters rather than arbitrary file groups.
Building the region map for $705$ nodes takes $\approx\!20$\,ms
($54.5$ builds/s); a single activation is $\approx\!20$\,ms.

\paragraph{Honest reading.} Absolute precision is low ($0.06$) because the v0.2
parser emits only containment and import edges, so $\Nk(t)$ is tiny (often $2$--$3$
nodes) while a region spans a whole file. The robust, real wins are
\emph{recall} ($0.97$ vs $0.35$) and \emph{token efficiency} ($\approx\!48\%$):
a region reliably contains a task's causal neighbourhood and pays fewer tokens to
do so. Richer semantic edges (call graph / data flow, roadmap item P1.3) would
sharpen $\Nk$ and is the main lever to raise precision; we flag this rather than
hide it.

\section{Hypothesis simulation under a stated oracle (LLM-free)}\label{sec:sim}

The core thesis --- \emph{does regional causal memory help an agent on long,
multi-file tasks?} --- ultimately needs LLM rollouts (\S\ref{sec:protocol}). But
its \emph{necessary condition} is testable now, without an LLM: \textbf{an agent
cannot solve a task whose required causal context is absent from its window}. We
measure, under an explicit oracle, whether each selection strategy \emph{places
the required causal context in the window}.

\paragraph{Setup.} We generate modular synthetic repositories: $S$ independent
subsystems, each a set of files linked by one cross-file causal chain plus
high-score \emph{decoy} symbols; there are no edges between subsystems, so each
subsystem is one bounded causal region. The causal structure is \emph{decoupled}
from the lexical structure --- a chain neighbour shares no identifier token with
the target, only an edge --- modelling a dependency that is causally essential yet
lexically dissimilar. A task of \emph{diameter} $d$ requires the $\pm d$ window
along its chain (up to $2d{+}1$ files); the budget holds $\approx$ one subsystem,
far less than the repository. The \textbf{oracle} is $R(t)\subseteq S$.

Crucially, retrieval pipelines locate code from a text \emph{query}, whereas an
OS-style memory \emph{anchors} on the workspace signal (the active file, a failing
test). We model both: each task carries a query (a token bag) and an anchor (the
active file node), and we run two scenarios --- \textbf{clean} (the query points
at the target) and \textbf{noisy} (a trap decoy in an unrelated subsystem
out-scores the target lexically). Six strategies share the budget:
\code{rag-dense} (top-$k$ lexical), \code{rag-hybrid} (lexical $+$ causal score),
\code{graphrag-1hop} (best hit $+$ one hop), \code{graphrag-bfs} (unbounded
expansion from the best hit), \code{ccos-from-query} (CCOS region of the best
\emph{lexical} hit --- an ablation), and \code{ccos-region} (CCOS region of the
\emph{anchor}). Seeded and deterministic; reproduce with \code{ccos experiment}.

\begin{table}[h]
\centering
\begin{tabular}{lcccccc}
\toprule
& \multicolumn{4}{c}{success at diameter $d$} & \multicolumn{2}{c}{overall} \\
\cmidrule(lr){2-5}\cmidrule(lr){6-7}
Strategy & $d{=}1$ & $d{=}2$ & $d{=}3$ & $d{=}4$ & succ. & cov. \\
\midrule
\multicolumn{7}{l}{\emph{Clean query (points at the target)}}\\
\code{rag-dense}      & 0.00 & 0.00 & 0.00 & 0.00 & 0.00 & 0.19 \\
\code{rag-hybrid}     & 1.00 & 0.00 & 0.00 & 0.00 & 0.23 & 0.65 \\
\code{graphrag-1hop}  & 1.00 & 0.00 & 0.00 & 0.00 & 0.23 & 0.58 \\
\code{graphrag-bfs}   & 1.00 & 1.00 & 1.00 & 1.00 & 1.00 & 1.00 \\
\code{ccos-from-query}& 1.00 & 1.00 & 1.00 & 1.00 & 1.00 & 1.00 \\
\code{ccos-region}    & 1.00 & 1.00 & 1.00 & 1.00 & 1.00 & 1.00 \\
\midrule
\multicolumn{7}{l}{\emph{Noisy query (a decoy out-scores the target lexically)}}\\
\code{rag-dense}      & 0.00 & 0.00 & 0.00 & 0.00 & 0.00 & 0.19 \\
\code{rag-hybrid}     & 1.00 & 0.00 & 0.00 & 0.00 & 0.23 & 0.65 \\
\code{graphrag-1hop}  & 0.00 & 0.00 & 0.00 & 0.00 & 0.00 & 0.19 \\
\code{graphrag-bfs}   & 0.00 & 0.00 & 0.00 & 0.00 & 0.00 & 0.00 \\
\code{ccos-from-query}& 0.00 & 0.00 & 0.00 & 0.00 & 0.00 & 0.00 \\
\textbf{\code{ccos-region}} & \textbf{1.00} & \textbf{1.00} & \textbf{1.00} & \textbf{1.00} & \textbf{1.00} & \textbf{1.00} \\
\bottomrule
\end{tabular}
\caption{Hypothesis simulation ($800$ tasks, seed $42$, budget $\approx$ one
subsystem). Success $=$ the required causal set $R(t)$ is inside the window. Under
a clean query every structure-aware method ties; under a misleading query only
the workspace-anchored region survives.}
\label{tab:sim}
\end{table}

\paragraph{Findings (and their honest limits).}
\textbf{(1)} Lexical retrieval (\code{rag-dense}) \emph{fails} on cross-file
causal tasks ($0\%$; coverage $0.19$): similarity cannot surface
causally-essential but lexically-dissimilar context --- the hypothesis's premise.
(\code{rag-hybrid}'s $d{=}1$ wins come from the \emph{causal score} it borrows,
not lexical similarity, which is why they persist under noise.)
\textbf{(2)} The value of structure-aware selection \emph{grows with diameter}:
one-hop expansion solves only $d{=}1$, while full causal paging
(\code{graphrag-bfs}, \code{ccos-region}) solves all $d$ --- the direction of H2.
\textbf{(3) The clean tie.} Under a clean query, \code{graphrag-bfs},
\code{ccos-from-query} and \code{ccos-region} all reach $1.00$: the lever is
causal \emph{structure}, \textbf{not CCOS per se}.
\textbf{(4) The noisy separation.} Under a \emph{misleading} query, every method
that locates code lexically collapses to $0\%$ --- including the strong
\code{graphrag-bfs} \emph{and} the \code{ccos-from-query} ablation (CCOS seeded on
the query). Only \code{ccos-region}, which anchors on the workspace signal rather
than the query, survives at $1.00$. The ablation isolates the differentiator: it
is the \emph{anchor source} (a structural workspace signal vs.\ a lexical query),
not the region machinery. This is the realistic regime --- a task description
rarely names the exact distant cause, and an OS-style memory that tracks the
active working set is robust where query-time retrieval is misled.
\textbf{(5) Assumptions, stated.} This credits CCOS with a reliable anchor (the
active file / failing test), which an OS-level memory has but a pure retrieval
pipeline does not; and it assumes \emph{modular} repositories with separable
regions (density $0.955$ on real code, \S\ref{sec:eval}) --- a monolithic giant
region exceeding the budget collapses the advantage (observed, reported).
\textbf{(6)} Throughout this is a \emph{simulation under a stated oracle}: it
tests the \emph{necessary} (retrieval) condition, not the \emph{sufficient}
(generation) one of \S\ref{sec:protocol}.

\section{Real-LLM evaluation: a first on-device measurement}\label{sec:protocol}

The simulation (\S\ref{sec:sim}) established the necessary condition. The
sufficient condition --- that an LLM agent then \emph{solves} more tasks with
regional memory than with chunk retrieval --- requires model rollouts. We
\textbf{implement the harness} (\code{ccos eval}) and report a \textbf{first
real-LLM run} against a local model.

\paragraph{Implemented harness.} Each task is a tiny multi-file project encoding
an \emph{arithmetic causal chain}: a base constant in one file is transformed
through a chain of one-line functions across separate files, and the question
``what integer does the last function return?'' is answerable \emph{only} by
reading the whole chain (the distant cause included). Grading is therefore
exact-match on an integer --- objective and automatable, with no code execution.
The same six strategies of \S\ref{sec:sim} assemble the context window under a
token budget (with the clean/noisy query split), the window is sent to a model
(any OpenAI-, Anthropic-Messages-, or Ollama-compatible endpoint), and we record
three metrics per diameter: \textbf{task success} (correct integer),
\textbf{oracle coverage} (required chain $\subseteq$ window --- model-independent),
and \textbf{symbol-hallucination} (the answer cites a function absent from the
project).

\paragraph{Setup.} We run \code{ccos eval} on-device on an NVIDIA Jetson AGX Thor
against \code{qwen2.5:7b-instruct} served by Ollama (temperature $0$), with $20$
tasks, seed $7$, a $600$-token budget, diameters $1$--$4$, in both the clean and
noisy regimes. Table~\ref{tab:real} reports the run.

\begin{table}[h]
\centering
\begin{tabular}{lcccccr}
\toprule
& \multicolumn{4}{c}{success at diameter $d$} & & \\
\cmidrule(lr){2-5}
Strategy & $d{=}1$ & $d{=}2$ & $d{=}3$ & $d{=}4$ & cov. & tok. \\
\midrule
\multicolumn{7}{l}{\emph{Clean query (names the target function)}}\\
\code{rag-dense}      & 0.12 & 0.00 & 0.00 & 0.00 & 0.00 & 519 \\
\code{rag-hybrid}     & 0.00 & 0.00 & 0.00 & 0.00 & 0.00 & 521 \\
\code{graphrag-1hop}  & 1.00 & 0.00 & 0.00 & 0.00 & 0.40 & 270 \\
\code{graphrag-bfs}   & 1.00 & 0.00 & 0.17 & 0.00 & 1.00 & 402 \\
\code{ccos-from-query}& 1.00 & 0.00 & 0.17 & 0.00 & 1.00 & 402 \\
\code{ccos-region}    & 1.00 & 0.00 & 0.17 & 0.00 & 1.00 & 402 \\
\midrule
\multicolumn{7}{l}{\emph{Noisy query (a decoy out-scores the target lexically)}}\\
\code{rag-dense}      & 0.00 & 0.00 & 0.00 & 0.00 & 0.00 & 531 \\
\code{rag-hybrid}     & 0.00 & 0.00 & 0.00 & 0.00 & 0.00 & 521 \\
\code{graphrag-1hop}  & 0.00 & 0.00 & 0.00 & 0.00 & 0.00 & 165 \\
\code{graphrag-bfs}   & 0.00 & 0.00 & 0.00 & 0.00 & 0.00 & 165 \\
\code{ccos-from-query}& 0.00 & 0.00 & 0.00 & 0.00 & 0.00 & 165 \\
\textbf{\code{ccos-region}} & \textbf{1.00} & \textbf{0.00} & \textbf{0.17} & \textbf{0.00} & \textbf{1.00} & \textbf{402} \\
\bottomrule
\end{tabular}
\caption{First real-LLM run: \code{qwen2.5:7b-instruct} via Ollama, on-device
(NVIDIA Jetson AGX Thor), $20$ tasks, seed $7$, budget $600$ tokens. ``cov.'' is
the model-independent oracle coverage; ``tok.'' the mean input tokens. Success is
an exact-match integer answer.}
\label{tab:real}
\end{table}

\paragraph{Findings.}
\textbf{(1) Coverage transfers from simulation to real text.} The
model-independent coverage column reproduces Table~\ref{tab:sim} on real file
content: lexical RAG covers $0$, the structure-aware methods cover $1.00$ under a
clean query, and under noise \emph{only} \code{ccos-region} holds $1.00$ --- the
necessary condition, confirmed outside the simulator.
\textbf{(2) Success is bounded by the model, not the context.} Where coverage is
$1.00$, the $7$B model still solves only $d{=}1$ ($1.00$) and part of $d{=}3$
($0.17$), missing $d{=}2,4$: with the whole chain in the window, the residual
failure is \emph{arithmetic reasoning} --- a model-capability ceiling, not a
selection failure. Crucially no method ever succeeds \emph{without} coverage
(success $\subseteq$ coverage), which is the claim the harness is built to test.
\textbf{(3) The noisy separation holds on a real LLM.} Under a misleading query
every query-driven method collapses to $0\%$ success --- including
\code{graphrag-bfs} and the \code{ccos-from-query} ablation --- while
\code{ccos-region}, anchored on the workspace signal, is the \emph{only} strategy
with non-zero success ($d{=}1$ at $1.00$). The query-driven methods even look
\emph{cheaper} under noise ($165$ vs $402$ tokens) precisely because they
confidently retrieve the wrong, smaller file; \code{ccos-region} spends its budget
on the causally-correct chain.
\textbf{(4) Honest limits.} $20$ tasks ($\approx 5$ per diameter) is a small
sample, and a $7$B model floors the sufficient condition. A frontier model
(\code{deepseek-v4-pro}, in progress) or a $70$B local model is expected to lift
success wherever coverage is $1.00$, \emph{sharpening} --- not changing --- the
separation, since coverage already fixes the ceiling.

\paragraph{Toward external benchmarks.} The arithmetic-chain suite isolates the
selection question; the decisive external tests are SWE-bench~\cite{jimenez2023swebench}
issue resolution (a patch passes the hidden tests) and a controlled
\emph{multi-file bug} suite where the fault site, its cause and its blast radius
lie in distinct files. Baselines share the base LLM and token budget: classical
RAG~\cite{lewis2020rag} (top-$k$ cosine over chunks), GraphRAG~\cite{edge2024graphrag}
(community summaries over a code graph), MemGPT~\cite{packer2023memgpt} (OS-style
paged memory), a LangGraph~\cite{langgraph} agent with a vector store, and CCOS
regions. Metrics: resolved-rate, input tokens to success, and a citation-grounding
hallucination rate checked against the ground-truth graph. We have begun this on
real history: \code{scripts/causal\_validation} mines fix commits, checks out the
pre-fix tree, injects the fault, and scores $R_{\mathrm{cov}} = |F_{\mathrm{target}}
\cap \mathrm{WorkingSet}_K| / |F_{\mathrm{target}}|$ --- the fraction of the files
a fix touched that the bounded working set recovers. On this repository's history
the first run exposed a limitation and its fix, both measured: with
downstream-only failure propagation $R_{\mathrm{cov}}$ was flat at $0.33$ (only the
seed file recovered, as co-changed files are \emph{upstream} importers reached
only through dependency hubs). Resolving intra-crate imports into file$\to$file
edges and propagating failure \emph{bidirectionally} fixes this. Across three
mature crates --- \code{fd}, \code{bat} and \code{hyperfine}, $70$ mined fix
commits --- the effect is consistent: at a sufficient budget ($K{\ge}50$) the two
changes lift $R_{\mathrm{cov}}$ to $0.85$--$1.0$ (from $0.50$--$0.84$
downstream-only), while \emph{diluting} to $0.19$--$0.28$ at a tight $K{=}20$
budget. \textbf{But against the obvious baseline, this is not a win.} Running a
classical lexical RAG (TF-IDF cosine over file text, queried by the fault file) at
the \emph{same file budget}, $R_{\mathrm{cov}}$ as CCOS\,/\,RAG is $0.92/0.94$,
$1.00/0.98$, $0.87/0.92$ at $K{=}50$ and ties at $K{=}100$, while at $K{=}20$ RAG
is clearly ahead ($0.20/0.56$ on \code{bat}, $0.20/0.73$ on \code{hyperfine}).
Causal selection therefore has \emph{no net coverage advantage} over lexical
similarity here, and is worse at a tight budget: on real bugs a fix's files are
lexically similar to one another, so TF-IDF recovers them too --- the §\ref{sec:sim}
premise of a cause lexically dissimilar from its symptom does not reproduce on
these repositories. The honest reading is that the \emph{necessary} condition
holds for the large majority but is \emph{not} CCOS-specific; a real advantage
would have to come from the regimes this setup does not test --- a degraded/absent
query (§\ref{sec:sim}), a symptom (failing test) lexically far from its cause
(needing the real test-driven seed, not the highest-degree heuristic), or the
\emph{sufficient} condition (Phase 4). Multi-crate workspaces, now linkable, also
remain to be measured at scale.

\paragraph{Sufficient condition (Phase 4): resolution ties, efficiency does not.}
We run the generation half on real single-file fixes from \code{fd}: an agent is
given an equal-budget context built two ways --- CCOS's causal region vs top-$k$
lexical RAG --- asked to rewrite the buggy file, and graded by whether
\code{cargo test} passes, with a \emph{compiler-in-the-loop} retry (on failure, a
\emph{context page fault} parses the error with the trace layer, injects pressure on
the faulting files, and re-prompts with a refreshed window). With a weak $7$B model
and one shot, both resolve $0/15$; with \code{qwen3-coder-30B} and three attempts
the loop lifts both to $2/10$ ($20\%$) --- the page-fault feedback, not the
retriever, unlocks resolution, and CCOS shows \textbf{no resolution advantage} over
RAG, consistent with the retrieval result. The \emph{efficiency} separates them: at
equal resolution CCOS's auto-bounded region averages $776$ context tokens against
the lexical retriever's budget-filling $5366$ --- a $\mathbf{6.9\times}$ reduction.
The same holds at scale without a model: across $51$ single-file fix scenarios from
\code{fd}, \code{bat} and \code{hyperfine}, CCOS assembles $700$--$1600$ context
tokens against RAG's $\approx\!6000$, a $\mathbf{4.1}$--$\mathbf{9.1\times}$
reduction. This is the one axis on which CCOS dominates: not \emph{what} it
retrieves (a baseline matches that) but \emph{how little} it needs, because the
causal region stops at the working set instead of padding a top-$k$ to budget. We
state the caveats: the \emph{resolution} sample is tiny ($n{=}10$, two passes), and
the baseline fills the budget by construction (a carefully tuned-$k$ RAG would also
be sparser). The defensible claim is narrower but real: CCOS \emph{self-calibrates}
--- it bounds itself at the causal region with no $k$ or budget to tune --- so it
never wastes the window, which is precisely the point of demand paging.

\paragraph{Hypotheses.}
\textbf{H1} (efficiency): CCOS reaches equal task success at fewer input tokens
than RAG/GraphRAG, because a region covers the causal neighbourhood with higher
recall (Table~\ref{tab:locality}).
\textbf{H2} (long-horizon success): on multi-file tasks CCOS's resolved-rate
exceeds chunk RAG by a margin that grows with the causal diameter of the task.
\textbf{H3} (grounding): CCOS lowers the symbol-hallucination rate, because the
admitted window is a connected sub-graph of \emph{real} nodes.

\paragraph{Threats to validity.} Region quality is bounded by edge quality
(\S\ref{sec:eval}); results would conflate the region \emph{policy} with the
underlying \emph{graph}. Ablations must therefore fix the graph and vary only the
selection strategy, and report per-task causal diameter so that gains are
attributed to coherence, not to retrieving more text. The \code{ccos-from-query}
ablation in Table~\ref{tab:real} is exactly this control --- same graph and
budget, only the anchor swapped --- and it collapses with the lexical baselines
under noise. A positive result on H1--H3 would constitute the research
contribution; a null result would still validate the deterministic, auditable
infrastructure.

\section{Limitations}

We are deliberately explicit about what is \emph{not} shown.
\textbf{(1) No coverage advantage over RAG on real bugs.} The headline real-data
metric ($R_{\mathrm{cov}}$, \S\ref{sec:protocol}) ties a plain lexical TF-IDF
retriever at a sufficient budget and loses to it at a tight one; on real fix
commits the files a fix touches are lexically similar, so the simulation's premise
of a lexically-dissimilar cause does not hold. The strong absolute numbers are a
\emph{necessary} condition that a baseline also satisfies, not a CCOS win.
(A complementary measurement on the engine's own source isolates a different,
structurally-defined relation: a lexical TF-IDF retriever recovers only $49\%$
(recall@$10$) of the AST-resolved \emph{import} dependencies---pairs that routinely
share no vocabulary, dependency cosine $0.53$ vs $0.43$ random---so RAG's blind spot
is real, but it lies off the fix-\emph{cohesion} axis the headline metric measures,
where a fix's files do share vocabulary. The structural layer recovers those edges by
construction; the now-default AST is what makes that recovery complete.)
\textbf{(2) Necessary $\neq$ sufficient.} We never measure whether a region
\emph{helps an agent fix} a bug (Phase 4: a patch passing the hidden tests) ---
only whether the relevant files are retrievable.
\textbf{(3) Seed heuristic.} The harness injects the fault at the
highest-out-degree changed file, not at the genuinely failing test, so it does not
exercise the regime where a symptom is lexically far from its cause --- the regime
most favourable to causal selection.
\textbf{(4) Scale and statistics.} Three single-crate repositories, $n\!\approx\!20$
each; differences are reported with their standard deviation but not cross-validated.
\textbf{(5) Engine.} The parser now \emph{defaults} to a real \texttt{syn} AST ---
measured $36.5\%$ more accurate than the former line heuristic on the engine's own
source (two-thirds vs.\ full import recall), the heuristic kept only as a fallback for
non-Rust input --- and ingestion was hardened to reconstruct bit-identically across
processes (sorted import resolution; order-invariant centrality). Regions remain
file-granular and merge only on explicit cross-file edges; scoring weights are
hand-tuned, though off-by-default eigenvector-centrality and node-lifecycle
(\texttt{Stable}/\texttt{Working}/\texttt{Orphan}) refinements are now available. None of
these affect the \emph{proven} properties (partition, determinism, replay) or the
measured locality --- but the empirical case for a downstream agent benefit is, at this
point, \emph{unproven}.

\section{Conclusion}

We set out to show that organising an agent's memory by \emph{causal regions}
retrieves long-horizon context better than chunk retrieval, gave the construction
a precise and deterministic definition, and proved it reconstructs bit-for-bit on
replay. We then tested the retrieval claim against the obvious baseline on real
data --- and it did not hold: across $70$ real bug-fix commits a plain lexical
TF-IDF retriever ties causal selection and beats it at a tight budget, and a
crash-trace pivot loses to RAG-over-the-error-message, because on real code a
fix's files and its error messages share vocabulary. We report the negative result
rather than bury it. What survives is not a better retriever but a different kind
of object: a \emph{deterministic, replayable, auditable} working memory in which an
agent's exact context state can be rewound and replayed under different
parameters --- time-travel debugging that a probabilistic retrieval stack cannot
offer. Whether that auditability, or a structured-context advantage at
\emph{generation} time (Phase 4), yields a measurable downstream gain remains open.
We release the engine, the honest validation harness and this paper so the
distinction between what is proven, what is measured, and what is refuted can be
checked rather than asserted.

\begin{thebibliography}{99}
\bibitem{lewis2020rag} P.~Lewis et al. Retrieval-Augmented Generation for
Knowledge-Intensive NLP Tasks. \emph{NeurIPS}, 2020. arXiv:2005.11401.
\bibitem{asai2023selfrag} A.~Asai et al. Self-RAG: Learning to Retrieve, Generate
and Critique through Self-Reflection. \emph{ICLR}, 2024. arXiv:2310.11511.
\bibitem{edge2024graphrag} D.~Edge et al. From Local to Global: A Graph RAG
Approach to Query-Focused Summarization. 2024. arXiv:2404.16130.
\bibitem{yamaguchi2014cpg} F.~Yamaguchi et al. Modeling and Discovering
Vulnerabilities with Code Property Graphs. \emph{IEEE S\&P}, 2014.
\bibitem{packer2023memgpt} C.~Packer et al. MemGPT: Towards LLMs as Operating
Systems. 2023. arXiv:2310.08560.
\bibitem{park2023generative} J.~S.~Park et al. Generative Agents: Interactive
Simulacra of Human Behavior. \emph{UIST}, 2023. arXiv:2304.03442.
\bibitem{langgraph} LangChain. LangGraph: Building Stateful, Multi-Actor
Applications with LLMs. Software framework, 2023--2024.
\url{https://github.com/langchain-ai/langgraph}.
\bibitem{denning1968working} P.~J.~Denning. The Working Set Model for Program
Behavior. \emph{Communications of the ACM}, 11(5), 1968.
\bibitem{liu2023lost} N.~F.~Liu et al. Lost in the Middle: How Language Models Use
Long Contexts. \emph{TACL}, 2024. arXiv:2307.03172.
\bibitem{kwon2023paged} W.~Kwon et al. Efficient Memory Management for Large
Language Model Serving with PagedAttention. \emph{SOSP}, 2023. arXiv:2309.06180.
\bibitem{haber1991timestamp} S.~Haber and W.~S.~Stornetta. How to Time-Stamp a
Digital Document. \emph{Journal of Cryptology}, 3(2), 1991.
\bibitem{cormack2009reciprocal} G.~V.~Cormack, C.~L.~A.~Clarke, and S.~B\"uttcher.
Reciprocal Rank Fusion Outperforms Condorcet and Individual Rank Learning Methods.
\emph{SIGIR}, 2009.
\bibitem{jimenez2023swebench} C.~E.~Jimenez et al. SWE-bench: Can Language Models
Resolve Real-World GitHub Issues? \emph{ICLR}, 2024. arXiv:2310.06770.
\end{thebibliography}

\end{document}
